% Section 4.1, from lectures/L04/appendix/a_adc.md.

\section{Analog-to-Digital Conversion (ADC)}
\label{c4:sec:adc}\label{c4:app:a}

\subsection{Why ADC}
\label{c4:a:1}
\begin{itemize}
  \item Everything so far has been binary: a pin is high or low. Real-world signals, like a
    potentiometer, a temperature sensor, or a microphone, vary continuously. An ADC
    (Analog-to-Digital Converter) samples that continuous voltage and returns a discrete number
    representing it.
  \item The ATmega328P has a 10-bit ADC: the result is an integer between \code{0} and \code{1023}
    ($2^{10} - 1$), representing an input voltage between \code{0 V} and whatever reference voltage
    you configure (normally the chip's own supply, \code{AVCC}, 5\,V on the Uno/Nano;
    \secref{c4:a:2} shows how to select it).
\end{itemize}
The mapping is linear:

\begin{console}
result = 1023 * (V_in / V_ref)
\end{console}

\noindent So \code{V_in = 0 V} $\rightarrow$ \code{0}, \code{V_in = V_ref} $\rightarrow$
\code{1023}, \code{V_in = V_ref / 2} $\rightarrow$ \code{~512}. More bits would mean finer
resolution for the same voltage range; 10 bits is enough for most sensor work, which is why the
chip stops there.

The datasheet writes the same mapping with \code{1024} instead of \code{1023}, counting full scale
as \code{V_ref} itself, so that the top code, \code{1023}, means one step below \code{V_ref}. The
two formulas never differ by more than one step (about 5\,mV at a 5\,V reference); this book uses
\code{1023} throughout, so that \code{1023} converts back to exactly \code{V_ref}.

There's a flip side to this same problem worth previewing: this chip's outputs are just as binary
as its inputs, a pin is either driving high or low, nothing in between. So how do you make an LED
look dimmer instead of just on or off? By toggling it faster than the eye can follow and varying
how much of each cycle it spends on versus off. That ratio, the fraction of one cycle spent ``on,''
is called the \textbf{duty cycle}; the \textbf{on-time} and \textbf{off-time} are just that
fraction (and its complement) converted into actual milliseconds for a chosen cycle length.
\Secref{c4:a:4} combines both ideas: an ADC reading drives the duty cycle, so turning a
potentiometer changes how bright an LED appears.

\subsection{Which pin, which registers}
\label{c4:a:2}
The chip exposes 8 analog input channels, \code{ADC0}--\code{ADC7} (plus a 9th, \code{ADC8}, wired
to an internal temperature sensor). On the Uno/Nano, \code{ADC0}--\code{ADC5} map to
\code{A0}--\code{A5}; you select a channel by writing its number into the low 4 bits of
\code{ADMUX}.

Three registers matter:
\begin{enumerate}
  \item \textbf{\code{ADMUX}} selects the input channel (bits 0--3) and the voltage reference
    (bits 6--7). Setting \code{REFS0} (and leaving \code{REFS1} clear) selects the supply voltage
    \code{AVCC} as the reference, the normal choice. Leave both clear, as they are after a reset,
    and the reference is whatever is connected to the \code{AREF} pin instead, which on an
    Uno/Nano is nothing at all:

\begin{ccode}
// Use AVCC as voltage reference, select ADC channel (0 for A0, 1 for A1 etc.).
ADMUX = (1U << REFS0) | channel;
\end{ccode}
  \item \textbf{\code{ADCSRA}} enables the ADC, starts a conversion, and sets its clock prescaler.
    The ADC needs a clock well below the CPU's; under $\sim$200\,kHz for full accuracy. At
    16\,MHz, a \textdiv128 prescaler gets you to 125\,kHz:

\begin{ccode}
// Enable the ADC (ADEN), start a conversion (ADSC), set prescaler to 128 (ADPSx).
ADCSRA = (1U << ADEN) | (1U << ADSC) | (1U << ADPS0) | (1U << ADPS1) | (1U << ADPS2);
\end{ccode}

    A conversion isn't instantaneous; poll the \code{ADIF} flag (or use an interrupt) until it's
    set, then clear it by writing a \code{1} back to it:

\begin{ccode}
while (0U == (ADCSRA & (1U << ADIF))) {} // Wait for the conversion.
ADCSRA = (1U << ADIF);                   // Clear the flag for next time.
\end{ccode}

    \noindent That plain \code{=} writes every other bit of \code{ADCSRA} too, \code{ADEN} included,
    so it also switches the ADC off until the next read turns it back on. That's harmless in
    \secref{c4:a:3}'s \code{adc_read()}, which rewrites the whole register at the start of every
    read, but it makes every conversion a first conversion after enabling, 25 ADC clock cycles
    (200\,\textmu s at 125\,kHz) instead of 13.
  \item \textbf{\code{ADC}} holds the result. In C this macro already combines the two underlying
    8-bit registers (\code{ADCH}/\code{ADCL}) into the full 10-bit value; you don't assemble it by
    hand. Without it, you'd have to do it yourself, and the read order matters:

\begin{ccode}
#define BYTE_SHIFT 8U // Parameter to shift an entire byte.

const uint8_t low = ADCL;  // Read ADCL first: locks ADCH until it's read.
const uint8_t high = ADCH; // Same conversion's result; safe to read now.
const uint16_t result = ((uint16_t)(high) << BYTE_SHIFT) | (uint16_t)(low);
\end{ccode}
\end{enumerate}

\note
\begin{itemize}
  \item Reading \code{ADCH} first risks tearing the result: if a new conversion finishes between
    the two reads, the two halves can end up belonging to different conversions.
  \item Reading \code{ADCL} first locks both registers against further updates until \code{ADCH}
    is read too, so the pair always comes from the same conversion; the \code{ADC} macro already
    does this for you.
  \item Both operands are cast to \code{uint16_t} before combining. The cast on \code{high} matters
    for correctness in general, not just for this specific value: \code{high} is a \code{uint8_t},
    and shifting it left by 8 without first widening it relies on C's automatic promotion to
    \code{int}.
  \item On AVR, \code{int} is 16 bits, so shifting a \emph{full-range} byte left by 8 (e.g.\
    \code{(uint8_t)0xFF << 8U}) can produce a value that no longer fits in a 16-bit signed
    \code{int}, and left-shifting past what the type can hold is undefined behavior in C.
  \item \code{ADCH} never actually uses more than its bottom 2 bits, so this particular shift is
    safe either way, but casting to \code{uint16_t} first is the habit that keeps it safe
    unconditionally: the shift then happens in an unsigned 16-bit type, where the result always
    wraps predictably instead of risking undefined behavior.
  \item The cast on \code{low} isn't load-bearing the same way (\code{uint8_t} always fits safely
    in the promoted type), but keeps both operands of \code{|} explicitly \code{uint16_t} instead
    of leaning on C's implicit-conversion rules to sort it out.
\end{itemize}

\subsection{A read function}
\label{c4:a:3}

\begin{ccode}
/**
 * @brief Read and return a 10-bit ADC conversion result.
 *
 * @param[in] channel ADC channel to read.
 *
 * @return Conversion result (0-1023).
 */
static uint16_t adc_read(const uint8_t channel)
{
    // Use AVCC as voltage reference and select ADC channel to read.
    ADMUX = (1U << REFS0) | channel;

    // Perform conversion, wait for conversion to complete.
    ADCSRA = (1U << ADEN) | (1U << ADSC) | (1U << ADPS0) | (1U << ADPS1) | (1U << ADPS2);
    while (0U == (ADCSRA & (1U << ADIF))) {}

    // Clear the ADC interrupt flag for next conversion, then return the result.
    ADCSRA = (1U << ADIF);
    return ADC;
}
\end{ccode}

\subsection{Example: potentiometer-controlled LED brightness (software PWM)}
\label{c4:a:4}
There's no true analog output on this chip (it has no digital-to-analog converter), but as
\secref{c4:a:1} previewed, you can fake brightness control in software instead: flick an LED on and
off fast enough that the eye averages it into perceived brightness, driven by duty cycle.

Potentiometer on A0, LED on D9, toggle button on D13 (enables/disables reading; the LED stays off
while disabled).

\bookimage{lectures/L04/appendix/images/circuit_pwm_with_pot.png}{A potentiometer on A0, an LED on
D9 and a button on D13.}{c4:fig:pot}

\begin{ccode}
#define F_CPU 16000000UL // CPU frequency in Hz.

#include <stdbool.h>
#include <stdint.h>

#include <avr/interrupt.h>
#include <avr/io.h>
#include <util/delay.h>

/** GPIO pins. */
#define LED1 1U // D9  -> PORTB1.
#define BTN1 5U // D13 -> PORTB5.
#define POT1 0U // A0  -> ADC0.

/** GPIO operations. */
#define LED1_ON PORTB |= (1U << LED1)      // Enable LED1.
#define LED1_OFF PORTB &= ~(1U << LED1)    // Disable LED1.
#define BTN1_PRESSED (PINB & (1U << BTN1)) // High if BTN1 is pressed, low otherwise.

/** Time parameters. */
#define PWM_PERIOD_MS 10U // PWM period in ms.
#define TICK_1MS 1U       // Generate a 1 ms tick in delay_ms().
#define ADC_MAX 1023U     // Maximum ADC conversion value.

/** ADC state (true = enabled, false = disabled). */
static volatile bool adc_state = false;

/**
 * @brief Set up system.
 */
static void setup(void)
{
    // Configure LED1 as output.
    DDRB = (1U << LED1);

    // Configure BTN1 as input with its internal pull-up enabled.
    PORTB = (1U << BTN1);

    // Enable pin change interrupts for BTN1.
    PCICR  = (1U << PCIE0);
    PCMSK0 = (1U << BTN1);
    sei();
}

/**
 * @brief Generate delay.
 *
 * @param[in] ms Delay duration in ms.
 */
static void delay_ms(const uint16_t ms)
{
    for (uint16_t i = 0U; i < ms; ++i)
    {
        _delay_ms(TICK_1MS);
    }
}

/**
 * @brief Read and return a 10-bit ADC conversion result.
 *
 * @param[in] channel ADC channel to read.
 *
 * @return Conversion result (0-1023).
 */
static uint16_t adc_read(const uint8_t channel)
{
    // Use AVCC as voltage reference and select ADC channel to read.
    ADMUX = (1U << REFS0) | channel;

    // Perform conversion, wait for conversion to complete.
    ADCSRA = (1U << ADEN) | (1U << ADSC) | (1U << ADPS0) | (1U << ADPS1) | (1U << ADPS2);
    while (0U == (ADCSRA & (1U << ADIF))) {}

    // Clear the ADC interrupt flag for next conversion, then return the result.
    ADCSRA = (1U << ADIF);
    return ADC;
}

/**
 * @brief Run a single PWM cycle.
 */
static void run_pwm_cycle(void)
{
    // Check if the ADC is enabled, terminate the function if not.
    if (!adc_state)
    {
        LED1_OFF;
        return;
    }

    // Read the ADC and compute the duty cycle.
    const uint16_t adc_val = adc_read(POT1);
    const float duty_cycle = (float)(adc_val) / ADC_MAX;

    // Compute the PWM durations based on the duty cycle, round to the nearest integer.
    const uint8_t on_ms  = (uint8_t)(duty_cycle * PWM_PERIOD_MS + 0.5F);
    const uint8_t off_ms = PWM_PERIOD_MS - on_ms;

    // Enable LED1 during the active duration.
    LED1_ON;
    delay_ms(on_ms);

    // Disable LED1 during the inactive duration.
    LED1_OFF;
    delay_ms(off_ms);
}

/**
 * @brief Toggle the ADC state on button press.
 */
ISR(PCINT0_vect)
{
    if (BTN1_PRESSED) { adc_state = !adc_state; }
}

/**
 * @brief Application entry point.
 *
 * @return 0 on termination of the program (should never occur).
 */
int main(void)
{
    setup();

    while (1)
    {
        run_pwm_cycle();
    }
    return 0;
}
\end{ccode}

\noindent Turning the potentiometer changes the ADC reading, which changes the computed duty cycle,
which changes how bright the LED appears.

Note \code{delay_ms()} is built on \code{_delay_ms()}, so it has the same blocking limitation
flagged in \secref{c2:a:6}; this software-PWM approach ties up the CPU for the whole 10\,ms
period. Chapter~\ref{c5:ch}, on timers, shows the way out: a hardware timer that raises an
interrupt when a period has passed, so the CPU is free in between instead of spinning inside
\code{_delay_ms()}. (The timers can also generate PWM entirely in hardware, which this book doesn't
cover.)
